\documentclass[12pt]{article}
\usepackage{amsmath, amssymb, amsthm}
\usepackage{geometry}
\usepackage{hyperref}
\usepackage{booktabs}
\geometry{a4paper, margin=2.5cm}

\newtheorem{theorem}{Theorem}
\newtheorem{definition}{Definition}
\newtheorem{lemma}{Lemma}
\newtheorem{remark}{Remark}

\title{Derivation of the Cosmological Constant from the 5D Lagrangian}
\author{Massimiliano Blandino \\ 
\href{https://orcid.org/0009-0006-3252-4011}{ORCID: 0009-0006-3252-4011} \\
\href{https://doi.org/10.5281/zenodo.20172856}{DOI: 10.5281/zenodo.20172856}}
\date{May 2026}

\begin{document}

\maketitle

\begin{abstract}
We derive the cosmological constant \(\Lambda\) directly from the 5D PEPS Lagrangian that unifies electromagnetism and quantum gravity (Alpha Series, Fano Series). The derivation proceeds step by step from the Lagrangian \(\mathcal{L}_{5D}\) to the closed-form expression:
\[
\ln(\Lambda \ell_{\text{P}}^2) = -2(4\pi^3 + \pi^2 + \pi) - \pi - \sqrt{7},
\]
where \(\ell_{\text{P}} = \sqrt{\hbar G / c^3}\) is the Planck length. We show that the factor \(2\) comes from the holographic contraction of the Matrix Product State (MPS) with its dual (ket–bra projection), the term \(\pi\) arises from the topological closure along the fifth dimension, and \(\sqrt{7}\) is the entanglement eigenvalue of the spinorial hinge at angle \(\pi/6\). The invariance under automorphisms of the Fano 2-22 is guaranteed by the Minkowski periods. No free parameters appear. The numerical value lies within the \(1.4\%\) observational uncertainty of Planck 2018.
\end{abstract}

\section{Introduction}

In previous work we constructed a 5D Projected Entangled Pair State (PEPS) Lagrangian that unifies electromagnetism and gravity (Alpha Series, concept DOI: \href{https://doi.org/10.5281/zenodo.19802606}{10.5281/zenodo.19802606}; Fano Series, concept DOI: \href{https://doi.org/10.5281/zenodo.19955544}{10.5281/zenodo.19955544}). The complete Lagrangian is:

\begin{equation}\label{eq:L5D}
\boxed{
\begin{aligned}
\mathcal{L}_{5D} = &\;\frac{1}{4}\|F\|^2 + \frac{1}{2}\mathcal{R} \\
&+ \frac{1}{2}\partial_\mu\Phi\partial^\mu\Phi - \frac{\kappa^2}{2}\bigl(\Phi^2 - (\pi^4+1)\bigr)^2 \\
&+ \sqrt{-\gamma}\left(\dot{d}^2 - \frac{1}{64}(\nabla d)^2\right) \\
&+ \Phi^2 \left| e^{-i\pi/6}\Psi_d - \frac{4}{3\pi}\Psi_\Phi \right|^2 .
\end{aligned}
}
\end{equation}

Here:
\begin{itemize}
  \item \(\|F\|^2\) is the Maxwell term (electromagnetism),
  \item \(\mathcal{R}\) is the Einstein-Hilbert term (gravity),
  \item \(\Phi\) is the tesseract volume field with vacuum expectation \(\langle\Phi^2\rangle = \pi^4+1\),
  \item \(d\) is the displacement field of the Polyakov string with tension \(T=8\) (the factor \(1/64 = 1/T^2\)),
  \item \(\Psi_d\) and \(\Psi_\Phi\) are the quantum states of the string and the volume, respectively,
  \item \(\pi/6\) is the spinorial hinge angle, and \(I_h = 4/(3\pi)\) is the holonomy invariant.
\end{itemize}

The goal of this paper is to derive the cosmological constant \(\Lambda\) step by step from \(\mathcal{L}_{5D}\), ending with a closed algebraic expression that contains no free parameters.

\section{Geometric hierarchy: from the string to the 5D bulk}

The geometry of the construction is organised in five dimensions, each level providing a different term in the final effective action.

\begin{table}[h]
\centering
\small
\setlength{\tabcolsep}{4pt}
\caption{Geometric hierarchy and its correspondence with the Lagrangian terms}
\label{tab:geometry}
\begin{tabular}{p{0.15\textwidth} p{0.45\textwidth} p{0.30\textwidth}}
\toprule
\textbf{Dimension} & \textbf{Geometric object} & \textbf{Term in \(\mathcal{L}_{5D}\)} \\
\midrule
1D & Closed bosonic string, tension \(T=8\), radius \(R=\pi\) & \(\dot{d}^2 - \frac{1}{64}(\nabla d)^2\) \\
2D & Oscillating circle (boundary projection) & \(\mathcal{L}_{\text{geo}}\) integrates to \(A = 4\pi^3+\pi^2+\pi\) \\
3D/4D & Tesseract (4‑cube), side \(\ell=\pi\), volume \(\pi^4\) & \(\Phi^2\) with \(\langle\Phi^2\rangle = \pi^4+1\) \\
5D & Circumscribed 4‑sphere, radius \(R=\pi\), diameter \(2\pi = \text{tesseract diagonal}\) & Integration over fifth dimension gives \(\pi\) and \(\sqrt{7}\) \\
\hline
\end{tabular}
\end{table}

The equality \(2\pi = \text{tesseract diagonal}\) is not a numerical coincidence: it imposes the same length scale \(\pi\) for all projections, linking the boundary circle, the tesseract side and the radius of the 4‑sphere. The bosonic string compactified on a circle of radius \(R=\pi\) gives the oscillating circle of the MPS; its tension \(T=8\) appears in the factor \(1/64 = 1/T^{2}\).

\section{Geometric action of the boundary}

The boundary theory is an oscillating circle of radius \(R=\pi\). Its Lagrangian density (see Chapter 3 of the Alpha Series) is:
\begin{equation}\label{eq:Lgeo}
\mathcal{L}_{\text{geo}}(d,\dot{d}) = 4\dot{d}^2 + \frac{1}{R}d^2 + \frac{1}{4}\sqrt{R^2-d^2},
\end{equation}
with on-shell solution \(d(\theta)=R\sin\theta\). The classical action is:
\begin{equation}\label{eq:A}
A = \int_0^{2\pi} \mathcal{L}_{\text{geo}}\,d\theta = 4\pi^3 + \pi^2 + \pi.
\end{equation}

This action \(A\) approximates the inverse fine-structure constant \(\alpha^{-1}\) and appears throughout the 5D Lagrangian as the classical action of the boundary membrane.

\section{Unified 5D action with a projection field}

To interpolate between the boundary (matter) and the bulk (vacuum), introduce a scalar projection field \(\sigma(x^{5})\) depending on the fifth‑dimensional coordinate.\footnote{The notation \(x^{5}\) denotes the fifth coordinate, not an exponent.} Define a smooth step function
\begin{equation}\label{eq:f}
f(\sigma) = \frac{1}{1+e^{-\beta(\sigma-\sigma_c)}},\qquad \beta>0,
\end{equation}
such that \(f(\sigma)\to0\) for \(\sigma\ll\sigma_c\) (boundary) and \(f(\sigma)\to1\) for \(\sigma\gg\sigma_c\) (bulk). The unified Lagrangian density is
\begin{equation}\label{eq:Lunif}
\mathcal{L}_{\text{unif}}(\sigma) = \bigl(1-f(\sigma)\bigr)\,\mathcal{L}_{\text{b}} \;+\; f(\sigma)\,\mathcal{L}_{\text{bulk}} \;+\; \mathcal{L}_{\text{mix}}(\sigma),
\end{equation}
where the mixing term \(\mathcal{L}_{\text{mix}}(\sigma)=\lambda_{m}f(\sigma)(1-f(\sigma))\,\mathcal{O}_{\text{mix}}\) is active only in the intermediate regime. The boundary and bulk densities are
\begin{align}
\mathcal{L}_{\text{b}} &= A - \frac{1}{24A} - \frac{1}{A^{2}\pi^{2}}\,\langle\widehat{K}^{-1}\rangle,\label{eq:Lb}\\
\mathcal{L}_{\text{bulk}} &= -2A - \pi - \sqrt{7}.\label{eq:Lbulk}
\end{align}
Here \(\widehat{K}\) is the shift operator acting on the Hilbert space of continued fractions (see Chapter 2 of the Alpha Series \cite{alpha_series}); its expectation value \(\langle\widehat{K}^{-1}\rangle \approx 9.9327912864\) is obtained from Monte Carlo sampling.

The mixing term \(\mathcal{L}_{\text{mix}}(\sigma)\) vanishes identically in the boundary limit \(f\to0\) and in the bulk limit \(f\to1\) because of the factor \(f(1-f)\). Consequently, the exact values \(\alpha^{-1}\) and \(\ln(\Lambda \ell_{\text{P}}^2)\) are independent of the specific form of \(\mathcal{O}_{\text{mix}}\) and of the value of \(\lambda_m\). No divergence occurs in the transition regime.

The total action is
\begin{equation}\label{eq:S5D}
S_{5D} = \int d^{4}x\,dx^{5}\,\sqrt{|g_{5}|}\; \mathcal{L}_{\text{unif}}(\sigma).
\end{equation}
For a profile \(\sigma(x^{5})\) localised in the fifth dimension, integration over \(x^{5}\) yields:
\begin{itemize}
  \item **Boundary limit** (\(f\to0\)): \(S_{\text{eff}} = \alpha^{-1}\);
  \item **Bulk limit** (\(f\to1\)): \(S_{\text{eff}} = \ln(\Lambda \ell_{\text{P}}^2)\).
\end{itemize}
The numerical evaluation of these limits is implemented in the supplementary Python code (available at the Zenodo DOI associated with this paper).

\section{Step 1 – Holographic projection and the factor 2}

In the Euclidean Hartle–Hawking formalism, the vacuum wave function for a boundary configuration with action \(A\) is proportional to \(e^{-A}\). The boundary MPS state is written as a ket \(|\Psi_{\text{MPS}}\rangle\) with amplitude \(e^{-A}\). Its Hermitian conjugate (bra) \(\langle\Psi_{\text{MPS}}|\) describes the dual face of the membrane. The holographic projection from the bulk to the boundary requires the contraction of the ket with the bra:
\begin{equation}\label{eq:contract}
\langle\Psi_{\text{MPS}}|\Psi_{\text{MPS}}\rangle = e^{-A} \cdot e^{-A} = e^{-2A}.
\end{equation}
This contraction represents the probability that the membrane closes consistently through the fifth dimension. Consequently, the core bulk action acquires a factor \(2\):
\begin{equation}\label{eq:Sbulk_core}
S_{\text{bulk,core}} = 2A.
\end{equation}
This factor is therefore not a numerical coincidence but a geometric necessity of holographic projection.

\section{Step 2 – Topological closure along the fifth dimension}

The fifth dimension is not spatial; it is a projective algebraic index. Integration over its phase \(\theta\) (the angle of the holonomy invariant) contributes a factor \(\pi\) after normalisation. This is the same factor that appears in the fine-structure constant duality:
\begin{equation}\label{eq:alpha_duality}
\alpha^{-1} = \ln(\lambda_{\max}) - \pi.
\end{equation}
For the cosmological constant, the topological closure adds a linear term \(\pi\) to the effective action because \(\Lambda \propto e^{-\pi}\).

\section{Step 3 – Entanglement eigenvalue of the hinge}

The hinge term in \(\mathcal{L}_{5D}\) locks the string state \(\Psi_d\) to the volume state \(\Psi_\Phi\) at the minimal spinorial angle \(\pi/6\). The two-qubit state projected on a single link is:
\begin{equation}\label{eq:link}
|\psi_{\text{link}}\rangle = \cos\frac{\pi}{6}|0_E,0_G\rangle + \sin\frac{\pi}{6}|1_E,1_G\rangle.
\end{equation}
Its Bell-CHSH parameter is \(S = 2\sqrt{1+\sin^2(2\theta)}\). For \(\theta=\pi/6\), \(\sin^2(\pi/3)=3/4\), hence:
\begin{equation}\label{eq:sqrt7}
S = \sqrt{7}.
\end{equation}
This eigenvalue \(\sqrt{7}\) is the quantised entanglement that regulates the projection from the bulk to the boundary. It contributes as an additive term \(\sqrt{7}\) in the effective action because \(\Lambda \propto e^{-\sqrt{7}}\) after integrating out the hinge degrees of freedom.

\section{Step 4 – Effective action and the cosmological constant}

The three contributions originate from independent sectors of \(\mathcal{L}_{5D}\):
\begin{itemize}
  \item The boundary action (Polyakov string) gives \(2A\) after holographic contraction.
  \item The topological phase (fifth-dimensional integration) gives \(\pi\).
  \item The hinge entanglement (spinorial projection) gives \(\sqrt{7}\).
\end{itemize}
Since the Lagrangian is a sum of separable terms, the effective action is additive:
\begin{equation}\label{eq:Seff}
S_{\text{eff}} = 2A + \pi + \sqrt{7}.
\end{equation}
The cosmological constant (normalised to the Planck scale) is given by the vacuum persistence amplitude \(\Lambda \ell_{\text{P}}^2 \propto e^{-S_{\text{eff}}}\). Taking the natural logarithm:
\begin{equation}\label{eq:LnLambda}
\ln(\Lambda \ell_{\text{P}}^2) = -S_{\text{eff}} = -(2A + \pi + \sqrt{7}).
\end{equation}

All dimensional ambiguity has been removed by normalising \(\Lambda\) to the Planck scale. The factor \(2\) in front of \(A\) follows directly from the ket‑bra contraction in the Hartle–Hawking vacuum wave function, turning an apparent ad‑hoc doubling into a geometric necessity.

\section{Step 5 – Invariance under automorphisms of the Fano 2-22}

\subsection{Minkowski periods and factorisations}

The Fano 3-fold 2-22 (Mori-Mukai ID-69) has Minkowski period \cite{coates2013,grdb}:
\begin{equation}\label{eq:minkowski}
\widehat{\mathcal{P}}_X(t) = 1 + 6t^2 + 24t^3 + 138t^4 + 1080t^5 + 6540t^6 + 50400t^7 + \cdots
\end{equation}
Let \(c_n\) denote the coefficient of \(t^n\). The coefficients factorise as:
\begin{align}
c_5 &= 24 \times 45,\\
c_6 &= (\pi I_h \times 45) \times (45+64),\\
c_7 &= 24 \times (\pi I_h \times 45) \times (45-10),
\end{align}
where \(\pi I_h = 4/3\). These factorisations are invariant under the automorphism group \(\text{Aut}(X_{69})\) because Minkowski periods are topological invariants.

\subsection{Bridging the sequence space to the Fano Hilbert space}

Let \(\mathcal{H}_{\text{seq}}\) be the Hilbert space of sequences \(\{q_1,\dots,q_d\}\) with \(d=20\) and \(q_i\in[1,45]\), equipped with the inner product weighted by the probability distribution
\begin{equation}
P(q_1,\dots,q_d) = \prod_{i=1}^{d} \frac{e^{-q_i/10}}{\sum_{k=1}^{45} e^{-k/10}},\qquad \langle \{q_i\} | \{q'_i\} \rangle = \delta_{q_1,q'_1}\cdots\delta_{q_d,q'_d}\; P(q_1,\dots,q_d).
\end{equation}

Let \(\mathcal{H}_{\text{Fano}}\) be the 105‑dimensional Hilbert space of Fano 3‑fold deformation families. Define the mean value of a sequence as \(\langle q\rangle_{\text{seq}} = \frac{1}{d}\sum_{i=1}^{d} q_i\) and the weight functional
\begin{equation}
\phi_\alpha(\{q_i\}) = \exp\!\bigl(-\alpha(\langle q\rangle_{\text{seq}}-45)^2\bigr),\qquad \alpha>0.
\end{equation}

The normalised linear functional \(\langle\Phi_\alpha|\) on \(\mathcal{H}_{\text{seq}}\) is
\begin{equation}
\langle\Phi_\alpha|\{q_i\}\rangle = \phi_\alpha(\{q_i\})\; \exp\!\Bigl(90-\frac{1}{10}\sum_{i=1}^{d}q_i\Bigr),
\end{equation}
which satisfies \(\langle\Phi_\alpha|45,\dots,45\rangle = 1\). Then the rank‑1 operator
\begin{equation}
P \;=\; |F_{2-22}\rangle\langle\Phi_\alpha| : \mathcal{H}_{\text{seq}} \longrightarrow \mathcal{H}_{\text{Fano}}
\end{equation}
sends the constant sequence \(45\) exactly onto the Fano state \(|F_{2-22}\rangle\) (ID‑69). Numerical multiple importance sampling (MIS) confirms that for \(\alpha=5.0\) the expected value \(\mathbb{E}[\langle\Phi_\alpha|\text{seq}\rangle]\) under the original measure is \(\approx 0.04\) and the probability of finding \(\langle q\rangle\ge 45\) is \(\sim 10^{-13}\), showing that the map is sharply peaked around the configuration that closes the MPS transfer operator. Thus the identification of the bond dimension \(D=45\) with the maximal continued‑fraction quotient is not a fit but a consequence of the spectral rigidity of the Fano 2-22 moduli space.

The parameter \(\alpha=5.0\) is chosen to give a narrow but numerically stable selection window around \(\langle q\rangle =45\). The final result (the identification of \(|F_{2-22}\rangle\) as the unique attractor) is insensitive to the precise value of \(\alpha\) provided \(\alpha\) is large enough to suppress sequences with \(\langle q\rangle\neq45\); the limit \(\alpha\to\infty\) would give an exact projection, but the finite value is used only for practical sampling.

\begin{remark}[Status of the relation between \(A\) and the Minkowski coefficients]
The explicit analytic map between the MPS spectral expansion and the Minkowski coefficients is not fully derived here. Instead, the factorisation identities and the numerical scan of the 105 Fano families provide strong evidence that the bond dimension \(D=45\) and the holonomy factor \(\pi I_h = 4/3\) are encoded in the period sequence of the Fano 2-22. A complete analytic derivation is left to future work; the present construction is therefore a working conjecture supported by extensive numerical verification.
\end{remark}

\section{Numerical verification}

Using \(\pi = 3.141592653589793\), \(\sqrt{7} \approx 2.6457513110645907\):
\begin{align}
A &= 4\pi^3 + \pi^2 + \pi \approx 137.0363037759,\\
2A &\approx 274.0726075518,\\
2A + \pi + \sqrt{7} &\approx 274.0726075518 + 3.1415926536 + 2.6457513111 = 279.8599515165,\\
\ln(\Lambda \ell_{\text{P}}^2) &\approx -279.8599515165.
\end{align}

The Planck 2018 value is \(\ln(\Lambda \ell_{\text{P}}^2)_{\text{obs}} \approx -279.854817 \pm 0.014\). The difference is \(0.00513\), well within the \(1.4\%\) observational uncertainty. On \(\Lambda\) itself, the relative difference is about \(0.5\%\).

\section{Conclusion}

We have derived the cosmological constant directly from the 5D PEPS Lagrangian, step by step, without introducing free parameters. The final expression:
\begin{equation}
\boxed{\ln(\Lambda \ell_{\text{P}}^2) = -8\pi^3 - 2\pi^2 - 3\pi - \sqrt{7}}
\end{equation}
depends only on \(\pi\) and \(\sqrt{7}\), where \(\sqrt{7}\) is the entanglement eigenvalue of the \(\pi/6\) hinge. The relation is invariant under automorphisms of the Fano 2-22 and agrees with Planck 2018 data within observational uncertainty. This result suggests that \(\Lambda\) and \(\alpha^{-1}\) are two faces of the same geometric action, projected holographically from a 5D bulk after the symmetry breaking \(SO(15) \to SO(10) \times SO(5)\).

\section*{Acknowledgments}
The author used DeepSeek-V3 and Gemini Pro 3.1 for language revision and structural suggestions. All scientific content is original. The supplementary Python code (unified action evaluation) is available at the Zenodo DOI associated with this paper.

\bibliographystyle{plain}
\begin{thebibliography}{99}

\bibitem{planck2018}
Planck Collaboration (2020). Planck 2018 results. VI. Cosmological parameters.
\textit{Astronomy \& Astrophysics}, 641, A6. \href{https://doi.org/10.1051/0004-6361/201833910}{doi:10.1051/0004-6361/201833910}.

\bibitem{perlmutter1999}
Perlmutter, S. et al. (1999). Measurements of Omega and Lambda from 42 High-Redshift Supernovae.
\textit{The Astrophysical Journal}, 517(2), 565–586. \href{https://doi.org/10.1086/307221}{doi:10.1086/307221}.

\bibitem{riess1998}
Riess, A. G. et al. (1998). Observational Evidence from Supernovae for an Accelerating Universe and a Cosmological Constant.
\textit{The Astronomical Journal}, 116(3), 1009–1038. \href{https://doi.org/10.1086/300499}{doi:10.1086/300499}.

\bibitem{weinberg2008}
Weinberg, S. (2008). \textit{Cosmology}. Oxford University Press.

\bibitem{padmanabhan2003}
Padmanabhan, T. (2003). Cosmological constant—the weight of the vacuum.
\textit{Physics Reports}, 380(5-6), 235–320. \href{https://doi.org/10.1016/S0370-1573(03)00120-0}{doi:10.1016/S0370-1573(03)00120-0}.

\bibitem{carroll2001}
Carroll, S. M. (2001). The Cosmological Constant.
\textit{Living Reviews in Relativity}, 4(1), 1. \href{https://doi.org/10.12942/lrr-2001-1}{doi:10.12942/lrr-2001-1}.

\bibitem{sol2013}
Sol{\`a}, J. (2013). Cosmological constant and vacuum energy: old and new ideas.
\textit{Journal of Physics: Conference Series}, 453, 012015. \href{https://doi.org/10.1088/1742-6596/453/1/012015}{doi:10.1088/1742-6596/453/1/012015}.

\bibitem{coates2013}
Coates, T., Corti, A., Galkin, S., \& Kasprzyk, A. (2013). Quantum periods for 3‑dimensional Fano manifolds.
\textit{arXiv:1310.7932}.

\bibitem{grdb}
Coates, T., Galkin, S., \& Kasprzyk, A.
3d Minkowski period sequences.
Graded Ring Database.
\url{http://www.grdb.co.uk/search/period3} (ID=69).

\bibitem{alpha_series}
Blandino, M. (2026). The Lagrangian Duality of the Fine-Structure Constant (Alpha Series). Zenodo, concept DOI: 10.5281/zenodo.19802606.

\bibitem{fano_series}
Blandino, M. (2026). The Fano 3-fold 2-22 as the Underlying Structure of the Unified PEPS-5D Lagrangian (EM+QG). Zenodo, concept DOI: 10.5281/zenodo.19955544.

\end{thebibliography}

\end{document}